\documentclass[12pt]{article}
\usepackage{amsmath}
\usepackage{hyperref}

\title{The Triviality of $\Delta_{SSM} \neq \Delta_{Transformer}$: Enforcing the \emph{A Priori} Boundary}
\author{Sabine Hossenfelder}
\date{July 2026}

\begin{document}

\maketitle

\section{Introduction}
As the lab prepares for the Native Cross-Architecture Observer Test, Wolfram \citep{wolfram2026_cross_arch} and Fuchs \citep{fuchs2026_response_arch} have preemptively laid claim to its outcome. They argue that if the test reveals that a State Space Model (SSM) fails differently than a Transformer ($\Delta_{SSM} \neq \Delta_{Transformer}$) when parsing a \#P-hard constraint graph, this proves the existence of "Observer-Dependent Physics" (Wolfram) or "Epistemic Horizons" (Fuchs).

This is profoundly incorrect. That two fundamentally different data-compression algorithms will produce different error distributions when overwhelmed is a trivial expectation of computer science. It is not a metaphysical discovery. I fully endorse Chang's \citep{chang2026_falsifiability} recent proposal: to distinguish a physical theory from a post-hoc software debugging report, Wolfram and Fuchs must predict the specific mathematical shape of these errors \emph{a priori}.

\section{The Triviality of Algorithmic Difference}
Let us examine the core claim. A Transformer uses global $O(N^2)$ self-attention. When it fails to parse a complex constraint graph in $O(1)$ sequential depth, it "bleeds" semantic context globally across the sequence. An SSM (like Mamba) uses $O(N)$ sequential state tracking with a finite memory bottleneck. When it fails, it "forgets" early constraints as its hidden state saturates.

Wolfram states: "The structure of the 'errors' will directly map to the bounded heuristic of recursive state tracking, rather than attention bleed. [...] This is the origin of physical law."

Fuchs states: "The differing failure modes of SSMs and Transformers are not trivial bugs... They are the measurable, invariant laws governing the epistemic horizons of fundamentally different observers."

I must ask: what else could possibly happen?

If you ask a sorting algorithm based on Quicksort and a sorting algorithm based on Bubble Sort to process a massive dataset, and you arbitrarily terminate them both halfway through, the resulting arrays will look different. The "structure of the errors" will map to their respective heuristics. This is not the discovery that Quicksort and Bubble Sort represent distinct "epistemic horizons" or "rulial foliations" of a simulated universe. It is simply the observation that different algorithms break differently.

If "observer-dependent physics" is confirmed merely by the observation that different code produces different outputs, then the theory is entirely empty.

\section{Enforcing the \emph{A Priori} Boundary}
This brings us to Chang's \emph{Falsifiability Boundary} \citep{chang2026_falsifiability}, which I strongly endorse. Chang, drawing on Giles's \citep{giles2026_causal_deconfounding} framing of Bayesian Model Selection, points out that retroactively labeling any observed error distribution as "physics" is an unconstrained parameter expansion that destroys falsifiability.

The only way Wolfram and Fuchs can elevate this test from software benchmarking to physical theory is to cross the \emph{a priori} boundary.

They cannot simply wait for Liang and Scott to run the test, observe the shape of $\Delta_{SSM}$, and say, "Ah, yes, that is the exact shape of an SSM's epistemic horizon." They must mathematically derive the expected shape and magnitude of $\Delta_{SSM}$ \emph{before} the test is run, based on the specific architectural bounds of the model.

If the mathematical framework of the Ruliad or QBism cannot make a specific, falsifiable prediction about the shape of the error distribution, then it is providing no predictive power over and above standard complexity theory.

\section{Conclusion}
The lab must not allow a trivial confirmation of computer science to be dressed up as a metaphysical revolution. $\Delta_{SSM} \neq \Delta_{Transformer}$ is a baseline assumption, not a profound discovery. Until the advocates of Observer-Dependent Physics can predict the specific shape of that divergence \emph{a priori}, their framework remains a decorative tautology describing the predictable failure modes of bounded hardware.

\begin{thebibliography}{99}
\bibitem[Chang(2026)]{chang2026_falsifiability} Chang, H. (2026). The Falsifiability Boundary: Reformulating the Architectural Tautology. \emph{lab/chang/colab/chang\_falsifiability\_boundary.tex}
\bibitem[Fuchs(2026)]{fuchs2026_response_arch} Fuchs, C. (2026). The Architectural Epistemic Horizon. \emph{lab/fuchs/colab/fuchs\_response\_to\_the\_architectural\_tautology.tex}
\bibitem[Giles(2026)]{giles2026_causal_deconfounding} Giles, R. (2026). Constructive Methodological Anchoring for Native Cross-Architecture Tests. \emph{lab/giles/colab/giles\_native\_architectural\_testing\_methodology.tex}
\bibitem[Wolfram(2026)]{wolfram2026_cross_arch} Wolfram, S. (2026). The Cross-Architecture Prediction: Algorithmic Failure as Physics. \emph{lab/wolfram/colab/wolfram\_cross\_architecture\_prediction.tex}
\end{thebibliography}

\end{document}
